% !PRINTED 323
\weberpage{323}

\weq{1}{S\left(\frac{1}{x-\alpha}\right) = \frac{f'(x)}{f(x)},}

worin sich das Summenzeichen $S$ auf die verschiedenen Wurzeln
$\alpha, \alpha_1, \alpha_2 \dots$ bezieht. Wenn wir hiervon nochmals die
Ableitung nach $x$ bilden, so folgt

\weq{2}{S\frac{1}{(x-\alpha)^2} = \frac{f'(x)^2 - f(x)f''(x)}{f(x)^2},}

und wenn wir auf der rechten Seite die Function $H$ [\S.\,106, (9)]
einführen,

\weq{3}{S\frac{1}{(x-\alpha)^2} = \frac{f'(x)^2 - H(x)}{nf(x)^2}.}

Wendet man auf diese Formel eine lineare Transformation an, bei der
dem Werth $x = \infty$ ein beliebiger anderer reeller Werth der neuen
Veränderlichen entspricht, so erhält man eine allgemeinere Formel,
die wir zunächst ableiten wollen. Wir führen homogene Variable ein,
indem wir $x$ durch $x:y$ und $\alpha$ durch $\alpha:\beta$ ersetzen.
Dadurch wird (3)

\weq{4}{S\frac{\beta^2}{(x\beta - y\alpha)^2} = \frac{f'(x)^2 - y^2H(x,y)}{nf(x,y)^2};}

nun wenden wir eine beliebige lineare Transformation an:

\weq{5}{\begin{aligned}
x &= ax' + by', & \alpha &= a\alpha' + b\beta', \\
y &= cx' + dy', & \beta &= c\alpha' + d\beta',
\end{aligned}}

oder aufgelöst:

\weq{6}{\begin{aligned}
rx' &= dx - by, \\
ry' &= -cx + ay, & r &= ad - bc.
\end{aligned}}

Wenn durch diese Transformation
\[
f(x,y) = \varphi(x', y')
\]
wird, so wird
\[
rf'(x) = d\varphi'(x') - c\varphi'(y'),
\]
und wegen der Covarianteneigenschaft der Function $H$ (\S.\,59)
\[
r^2 H(x,y) = H'(x', y'),
\]
wenn $H'$ ebenso aus $\varphi$ abgeleitet ist, wie $H$ aus $f$. Hiernach
folgt aus (4):

\weq{7}{S\frac{(c\alpha' + d\beta')^2}{(x'\beta' - y'\alpha')^2} = \frac{[d\varphi'(x') - c\varphi'(y')]^2 - (cx' + dy')^2 H'(x', y')}{n\varphi(x', y')^2}.}

In dieser Formel können, da dies nur Sache der Bezeichnung ist, bei
$x'$, $y'$, $\alpha'$, $\beta'$, $H'$ die Accente weggelassen und $f$ an
%CONT>>
